\documentclass[10pt,a4paper]{amsart}
\usepackage[margin=19mm]{geometry}
\usepackage{amsmath,amssymb,mathtools}
\usepackage[scr=rsfs,cal=euler]{mathalfa}
\usepackage{xfrac}
\usepackage[unicode,hidelinks,bookmarks=false]{hyperref}
\newcommand{\ie}{i.e.\ }
\newcommand{\cf}{cf.\ }
\newcommand{\resp}{respectively\ }
\newcommand{\Subsection}{\subsection}
\providecommand{\nobreakdash}{\nobreak\hbox{-}}
\setlength{\emergencystretch}{2em}
\multlinegap0pt
\usepackage{siunitx}
\usepackage{siunitx}
\usepackage{siunitx}
\usepackage{siunitx}
\usepackage{mathtools}
\usepackage{amssymb}
\usepackage{braket}
\usepackage{siunitx}
\usepackage{cancel}
\usepackage{tikz}
\usepackage[shortlabels]{enumitem}
\usepackage{csquotes}
\newtheorem{lemma}{Lemma}
\newcommand{\C}{\numberset{C}}
\renewcommand{\Im}{\Impart}
\DeclareMathOperator{\Impart}{Im}
\DeclareMathOperator{\Mat}{Mat}
\newcommand{\R}{\numberset{R}}
\renewcommand{\Re}{\Realpart}
\DeclareMathOperator{\Realpart}{Re}
\newcommand{\numberset}{\mathbb}
\renewcommand{\phi}{\varphi}
\renewcommand{\vec}{\boldsymbol}
\title{Analytic structure of $q$-pseudoconcave subsets of continuous graphs}
\author{Filippo Valnegri}
\date{Source: arXiv:2603.04967v3 (CC BY 4.0)}
\begin{document}
\maketitle

\noindent\emph{Source and licence.} Analytic structure of $q$-pseudoconcave subsets of continuous graphs. Version arXiv:2603.04967v3 (CC BY 4.0), distributed by arXiv under the Creative Commons Attribution 4.0 International licence, http://creativecommons.org/licenses/by/4.0/, as recorded in the arXiv licence metadata for this exact version and shown on the abstract page https://arxiv.org/abs/2603.04967v3. Full licence notice: \texttt{licenses/arxiv-2603.04967-CC-BY-4.0.txt}.\par\smallskip

\noindent\textit{Editorial scope.} Counted unit: the article's lemma lemma:reduction (source0.tex lines 709--720). The article's complete original proof of it is delivered with it (source0.tex lines 721--778, one \texttt{proof} environment of 58 lines). No mathematics is written, repaired or re-derived: the statement and the proof are the article's own lines, in the article's own notation, and no retained line is substituted or re-flowed.  No further source-local context is needed: every cross-reference of every retained line resolves inside the two delivered spans.  Nothing was omitted from any retained span, so the package carries no omission record.  No retained line contains a citation, so the wrapper carries no bibliography and no bibliography entry.  No retained line contains a figure environment, an image-inclusion command or drawing code, so no image is delivered and the package declares no asset dependency.  Editorial formatting only, in addition: an AMS \texttt{amsart} preamble that restates the article's own declarations and macros used by the retained lines, namely the environments \texttt{lemma} on separate counters, together with the standard packages those lines need.  The retained environments print in this document as Lemma 1 in order of appearance, whereas the article prints the same items with the numbering of its own full document.  The complete native source of the article as distributed in the arXiv e-print for this exact version is bundled unchanged as \texttt{sources/195772/source0.tex}.
\begin{lemma}
\label{lemma:reduction}
    For $n\geq1$, let $D\subset\C^n\times\R$ be a domain and let $X_0\subset D$ be a closed subset. For $N>n$ and $p=N-n-1\geq0$, let $g=(g',g''):X_0\to\R\times\C^p$ be a continuous function and let
    \[
        Z_0\coloneqq\Gamma(g)=\bigl\{(z,w,\zeta)\in X_0\times i\R\times\C^p \mid \bigl(\Im(w),\zeta\bigr)=g\bigl(z,\Re(w)\bigr)\bigr\}
    \]
    be its graph in $\C^N=\C^n_z\times(\R\times i\R)_w\times\C^p_\zeta$. Choose $r\in\{0,\dots,p\}$ and let
    \[
        \pi:D\times i\R\times\C^p\to D\times i\R\times\C^r, \quad \pi(z,w,\zeta)=(z,w,\zeta_{\mu_1},\dots,\zeta_{\mu_r})
    \]    
    be the projection, for $1\leq\mu_1<\dots<\mu_r\leq p$. If $Z_0$ is a $q$-local maximum set, then $Z'_0\coloneqq\pi(Z_0)$ is a $q$-local maximum set.
\end{lemma}

\begin{proof}
    Since we want to prove a local property, we can assume, at most shrinking $D$ and applying a holomorphic change of coordinates, that $\mu_j=j$ for all $j=1,\dots,r$.

    By contradiction, suppose that there exists a compact $K\subset D\times i\R\times\C^r$, a smooth $q$-plurisub\-harmonic function $\phi$, defined on a neighbourhood $U\subset D\times i\R\times\C^r$ of $K$, and a point $\eta_0\in Z'_0\cap\mathring{K}$ such that
    \begin{equation}
    \label{eq:contradiction}
        \phi(\eta_0)=\max_{Z'_0\cap K}\phi>\max_{Z'_0\cap bK}\phi
    \end{equation}
    Let us define the function
    \[
        \tilde\phi\coloneqq\phi\circ\pi:\pi^{-1}(U)=U\times\C^{p-r}\to\R
    \]
    We want to prove the $\tilde\phi$ is a smooth $q$-plurisubharmonic function. Note that
    \begin{equation}
    \label{eq:projection}
        \tilde\phi(z,w,\zeta)=\phi(z,w,\zeta_1,\dots,\zeta_r)
    \end{equation}
    and thus $\tilde\phi$ is a smooth function. Let
    \[
        \operatorname{L}_{\tilde\phi}=\biggl(\frac{\partial^2\tilde\phi}{\partial\xi_i \partial\bar{\xi}_j}\biggr)_{ij}\in\Mat_{\mathscr{C}^\infty(U\times\C^{p-r})}(N\times N)
    \]
    be the Levi matrix of $\tilde\phi$. Using~\eqref{eq:projection}, we obtain that
    \begin{equation}
    \label{eq:matrix}
        \operatorname{L}_{\tilde\phi}=
        \begin{pmatrix}
            \operatorname{L}_\phi & \vec{0}_{(n+1+r)\times(p-r)} \\
            \vec{0}_{(p-r)\times(n+1+r)} & \vec{0}_{(p-r)\times(p-r)}
        \end{pmatrix}
    \end{equation}
    where $\operatorname{L}_\phi\in\Mat_{\mathscr{C}^\infty(U)}((n+1+r)\times(n+1+r))$ is the Levi matrix of $\phi$. If $\sigma(\operatorname{L}_\phi(\xi))=\{\alpha_1,\dots,\alpha_{n+1+r}\}$ is the spectrum of $\operatorname{L}_\phi(\xi)$ (with repetitions allowed), for a fixed point $\xi\in U$, then
    \[
        \sigma(\operatorname{L}_{\tilde\phi}(\xi))=\{\alpha_1,\dots,\alpha_{n+1+r},0,\dots,0\}
    \]
    by~\eqref{eq:matrix}. Since $\phi$ is $q$-plurisubharmonic, then $\sigma(\operatorname{L}_\phi(\xi))$ has at most $q$ negative elements and thus $\sigma(\operatorname{L}_{\tilde\phi}(\xi))$ has at most $q$ negative elements. Therefore $\tilde\phi$ is a smooth $q$-plurisubharmonic function.

    Note that
    \begin{align*}
        Z_0 &= \bigl\{\bigl(x,ig'(x),g''(x)\bigr) \mid x\in X_0\bigr\} \\
        Z'_0 &= \bigl\{\bigl(x,ig'(x),g''_1(x),\dots,g''_r(x)\bigr) \mid x\in X_0\bigr\}
    \end{align*}
    Since $\eta_0\in Z'_0$, then $\eta_0=(x_0,ig'(x_0),g''_1(x_0),\dots,g''_r(x_0))$ for some $x_0\in X_0$ and thus
    \[
        \pi^{-1}(\eta_0)\cap Z_0=\{(x_0,ig'(x_0),g''(x_0)\}=\{(\eta_0,\chi_0)\}
    \]
    where $\chi_0\coloneqq(g''_{r+1}(x_0),\dots,g''_p(x_0))$. By~\eqref{eq:contradiction}, we obtain that
    \begin{equation}
    \label{eq:maxpoint}
        \tilde\phi(\eta_0,\chi_0)=\phi(\eta_0)>\phi(\eta)=\phi(\eta,\chi)
    \end{equation}
    for all $(\eta,\chi)\in\pi^{-1}(Z'_0\cap bK)$.

    Take $C\subset\C^N$ compact such that $(\eta_0,\chi_0)\in Z_0\cap\mathring{C}$, $\pi(C)=K$ and $Z_0\cap bC= Z_0\cap\pi^{-1}(bK)$. Since $Z_0\cap bC\subset\pi^{-1}(Z'_0\cap bK)$, then
    \[
        \tilde\phi(\eta_0,\chi_0)>\max_{Z_0\cap bC}\tilde\phi
    \]
    by~\eqref{eq:maxpoint}. This is a contradiction with the $q$-local maximum property of $Z_0$. Therefore $Z'_0$ is a $q$-local maximum set.
\end{proof}

\end{document}
